\documentclass{situationist-preprint}

\title{The Function of Boredom}
\msid{EXI-10-a}
\msfamily{EXI-10}
\msdiscipline{Meaning, Death and the Absurd}
\mstype{Principal manuscript}
\mskeywords{boredom; function; emotion; engagement; attention; Heidegger; Elpidorou; Svendsen}

\begin{document}
\maketitle

\begin{abstract}
What is boredom for? It is unpleasant, often treated as a vice or a symptom, and in the existentialist tradition as a revelation of the emptiness of things. We argue that boredom has a function: it signals a mismatch between one's capacities and one's activities, and its unpleasantness motivates a change of activity that tends to restore the match. We give a selected-effect account of function and argue that boredom has the features that such an account requires. It is a state with content, which represents the current situation as failing to engage one's attention and abilities; it is triggered by the conditions in which that representation is correct, namely under-challenge, over-challenge without traction, and lack of meaning; and it motivates search for alternatives, which in environments like those in which the capacity was shaped tends to improve the match. The account explains why boredom is aversive, why it is sensitive to meaning as well as to difficulty, and why chronic boredom is a malfunction or a correct signal that cannot be acted on. It leaves the Heideggerian thesis about profound boredom in place as a claim about a different, rarer state.
\end{abstract}

\section{The problem}

Boredom has had a bad reputation. In the tradition of the desert monks it appears as \emph{acedia}, a vice of spiritual sloth; for Schopenhauer it is evidence that existence has no positive value; for \citet{Svendsen2005} it is a characteristic affliction of modernity, the condition of a culture in which meaning has to be produced rather than inherited. \citet{Heidegger1995}, in lectures of 1929--30, distinguished being bored by something, becoming bored with something, and profound boredom, in which beings as a whole withdraw and Dasein is brought before its own temporality. On these views boredom is a fault, in the world or in us, and its question is what it reveals.

A different question has been raised by recent philosophy and psychology: what boredom does. \citet{Elpidorou2014, Elpidorou2018} argues that boredom is a regulatory state that informs us that our situation does not fit our desires and goals and motivates us to pursue a more satisfying one. \citet{BenchLench2013} propose that boredom encourages the pursuit of alternative goals when the current one has ceased to be rewarding. \citet{EastwoodEtAl2012} define boredom as the aversive experience of wanting, but being unable, to engage in satisfying activity, and locate its source in attention. \citet{WestgateWilson2018} distinguish attentional and meaning components: one is bored when one's attention cannot be held, and when the activity does not matter to one.

This manuscript defends the functional view and makes it precise. Our conclusion is that boredom signals a mismatch between one's capacities and one's activities, and has an adaptive function. The task is to say what \enquote{function} means here, to show that boredom meets the conditions, and to show what the account explains that its rivals do not.

\section{Framework}

\begin{definition}[Boredom]\label{d:bored}
\emph{Boredom} is an aversive affective state in which the subject experiences her current activity or situation as failing to engage her, together with a desire to be engaged otherwise.
\end{definition}

\begin{definition}[Mismatch]\label{d:mismatch}
There is a \emph{mismatch} between an agent's capacities and her activity when the activity does not occupy her available attention and abilities in a way she finds worth sustaining: because it demands too little of them (under-challenge), because it demands more than she can bring to bear so that she gains no traction (over-challenge), or because, whatever its demands, it serves nothing she cares about (meaninglessness).
\end{definition}

\begin{definition}[Function]\label{d:function}
A trait has the \emph{function} of producing an effect $E$ if trait-tokens of its type were retained, by natural selection or by learning, because they produced $E$ \citep{Millikan1984}.
\end{definition}

\begin{principle}[Emotional content]\label{p:content}
Affective states represent their circumstances as having evaluative properties, and are correct when the circumstances have them. This is the position of MND-17, applied to boredom: boredom represents the current situation as failing to engage.
\end{principle}

\begin{assumption}[Value of engagement]\label{a:value}
For creatures whose survival and flourishing depend on learning, exploration and the exercise of skill, time spent in activities that match capacities yields, on average, more learning, skill and goal attainment than time spent in activities that do not.
\end{assumption}

The three clauses of Definition~\ref{d:mismatch} correspond to the three triggers documented in the empirical literature: monotony, tasks that are too hard to make progress on, and tasks perceived as pointless \citep{WestgateWilson2018}. The meaning clause is why the mismatch is between capacities and \emph{activities} rather than between capacities and task difficulty alone: the capacities in question include the capacity to care, and an activity that cannot be connected to anything one cares about leaves it idle. Assumption~\ref{a:value} is the premise that \citet{Csikszentmihalyi1990} documented in the opposite direction: absorption and learning are maximal where challenge matches skill.

\section{The argument}

\begin{lemma}\label{l:signal}
Boredom represents a mismatch, and is correct when there is one.
\end{lemma}

\begin{proof}
By Principle~\ref{p:content} and Definition~\ref{d:bored}, boredom represents the current activity as failing to engage the subject. An activity fails to engage an agent who is not otherwise impaired just when it does not occupy her attention and abilities in a way she finds worth sustaining, which is a mismatch (Definition~\ref{d:mismatch}). So boredom represents a mismatch, and its correctness condition is that there be one.
\end{proof}

\begin{lemma}\label{l:motiv}
Boredom motivates the search for an alternative activity, and the search tends to reduce mismatch.
\end{lemma}

\begin{proof}
Boredom is aversive and includes a desire to be engaged otherwise (Definition~\ref{d:bored}). An aversive state motivates action that ends it, and the action that ends boredom is a change of activity, or a change in how one engages with the present one, that removes the mismatch. Experimental evidence confirms the motivational profile: bored subjects prefer novel and even unpleasant stimulation to continued monotony \citep{WilsonEtAl2014}. A search that terminates when engagement is restored selects, among available alternatives, those that engage.
\end{proof}

\begin{theorem}\label{t:main}
Boredom signals a mismatch between capacities and activities, and has the adaptive function of moving the agent toward activities that match her capacities.
\end{theorem}

\begin{proof}
The signaling claim is Lemma~\ref{l:signal}. By Lemma~\ref{l:motiv}, boredom produces the effect of moving the agent toward matching activities. By Assumption~\ref{a:value}, agents who are so moved spend more of their time in activities that yield learning, skill and goal attainment, and these contribute to survival and flourishing. A disposition to boredom was therefore advantageous to its bearers, and is retained, in the species and in individual development, because it produces that effect. By Definition~\ref{d:function}, moving the agent toward matching activities is its function, and the function is adaptive.
\end{proof}

The account explains two features of boredom. It explains why boredom is \emph{aversive}: a signal whose function is to end the current activity must motivate ending it. It explains why boredom is \emph{sensitive to meaning}: the mismatch it signals includes the idleness of one's capacity to care, so that a demanding but pointless task can bore as much as an easy one.

\begin{corollary}\label{c:chronic}
Chronic boredom is either a malfunction, when the signal is produced in the absence of mismatch, or a correct signal that cannot be acted on, when the agent is confined to mismatched activity.
\end{corollary}

\begin{proof}
A signal with the function of Theorem~\ref{t:main} fails to perform it in two ways: when it represents a mismatch that is not there, which by Lemma~\ref{l:signal} makes it incorrect, or when it is correct but the motivated search has no available alternatives to find. Persistent boredom is a failure of the function, so it is one of these.
\end{proof}

The corollary separates two conditions that the cultural critique of boredom runs together. Boredom proneness associated with depression or attentional impairment is plausibly of the first kind. The boredom of the prisoner, the assembly-line worker or the child in a class pitched far below her is of the second: it represents its situation correctly, and its persistence is a fact about the situation.

\section{Objections and replies}

\begin{objection}[Profound boredom]\label{o:heid}
Heidegger's profound boredom is not about any activity. In it, everything is indifferent, and no change of activity is called for or helps. A functional account cannot capture it \citep{Heidegger1995}.
\end{objection}

\begin{reply}
We agree that profound boredom, if it occurs as Heidegger describes it, is not the state of Definition~\ref{d:bored}. It has a different object: PHE-04 argues that it is intentional toward time itself, not toward an activity. The functional account concerns the ordinary boredom of situations and tasks, which is the boredom with an evident motivational role. The existence of a distinct and rarer state bearing the same name does not show that the ordinary one lacks a function.
\end{reply}

\begin{objection}[Boredom harms]\label{o:harm}
Boredom is associated with overeating, gambling, substance use and risk-taking. A state with these effects is not adaptive.
\end{objection}

\begin{reply}
The objection confuses the function of a state with the effects of the responses available to it. Boredom motivates a search for engagement; it does not select which alternatives are on offer. In environments rich in cheap stimulation, the search terminates in whatever restores arousal fastest, which need not be what restores the match. This is a mismatch between the environment and the conditions under which the disposition was retained, of the kind that explains other maladaptive effects of adaptive traits. The harmful responses are, moreover, short-lived remedies that leave the mismatch in place, which is why the boredom returns.
\end{reply}

\begin{objection}[Function is not value]\label{o:normative}
Even if boredom has a biological function, that does not show that we should heed it. Evolution does not settle what is good for us.
\end{objection}

\begin{reply}
The conclusion of this family is that boredom has an adaptive function, not that every episode should be obeyed. Its signal is defeasible: a bored student may have good reason to finish the dull proof. But the signal bears on a question that matters independently of its evolutionary history: whether one's capacities are being put to use. That is a question about the agent's good on the account of meaning in EXI-17, on which meaning requires active engagement with something worthwhile, and boredom is a reliable, if fallible, indication of the absence of engagement.
\end{reply}

\section{Conclusion}

Boredom is the felt registration of a mismatch between what one can bring to an activity and what the activity asks of one or offers to one, including whether it serves anything one cares about. It motivates a change that tends to restore the match, and because creatures who must learn and exercise skills do better when their capacities are engaged, the disposition to boredom has the adaptive function of moving them toward engagement. Chronic boredom is a malfunction of the signal or a correct signal that cannot be acted on. The account is consistent with EXI-03's thesis that an endless life would exhaust the desires that give it its point: in such a life, boredom would become a correct signal that nothing could answer.

\begin{thebibliography}{99}\small
\bibitem[Bench and Lench(2013)]{BenchLench2013} Bench, S.~W. and Lench, H.~C. (2013). On the function of boredom. \emph{Behavioral Sciences}, 3(3), 459--472.
\bibitem[Csikszentmihalyi(1990)]{Csikszentmihalyi1990} Csikszentmihalyi, M. (1990). \emph{Flow: The Psychology of Optimal Experience}. Harper \& Row.
\bibitem[Eastwood et~al.(2012)]{EastwoodEtAl2012} Eastwood, J.~D., Frischen, A., Fenske, M.~J. and Smilek, D. (2012). The unengaged mind: Defining boredom in terms of attention. \emph{Perspectives on Psychological Science}, 7(5), 482--495.
\bibitem[Elpidorou(2014)]{Elpidorou2014} Elpidorou, A. (2014). The bright side of boredom. \emph{Frontiers in Psychology}, 5, 1245.
\bibitem[Elpidorou(2018)]{Elpidorou2018} Elpidorou, A. (2018). The good of boredom. \emph{Philosophical Psychology}, 31(3), 323--351.
\bibitem[Heidegger(1995)]{Heidegger1995} Heidegger, M. (1995). \emph{The Fundamental Concepts of Metaphysics: World, Finitude, Solitude}. Trans. W. McNeill and N. Walker. Indiana University Press.
\bibitem[Millikan(1984)]{Millikan1984} Millikan, R.~G. (1984). \emph{Language, Thought, and Other Biological Categories}. MIT Press.
\bibitem[Svendsen(2005)]{Svendsen2005} Svendsen, L. (2005). \emph{A Philosophy of Boredom}. Trans. J. Irons. Reaktion Books.
\bibitem[Westgate and Wilson(2018)]{WestgateWilson2018} Westgate, E.~C. and Wilson, T.~D. (2018). Boring thoughts and bored minds: The MAC model of boredom and cognitive engagement. \emph{Psychological Review}, 125(5), 689--713.
\bibitem[Wilson et~al.(2014)]{WilsonEtAl2014} Wilson, T.~D., Reinhard, D.~A., Westgate, E.~C., Gilbert, D.~T., Ellerbeck, N., Hahn, C., Brown, C.~L. and Shaked, A. (2014). Just think: The challenges of the disengaged mind. \emph{Science}, 345(6192), 75--77.
\end{thebibliography}

\end{document}
